% Nouveau EGA V (cote 25): archive cote 25, site page 17 (PDF page 18).
% typescript. Draft compared: none.
% AI-transcribed by Claude Opus 5.5 (claude-opus-5-5, adaptive thinking, effort high), Message Batches API, batch ega25b (msgbatch_01A7n5YuXubUPNK1zsepi8RG), 2026-11-13; not yet checked by eye.
\scanpage{25/17}
\mspage{14}

(NB il y aurait lieu d'inclure cette description birationnelle au moins
en corollaire à 2.1. ). La démonstration de \over{\unsure{2}}{3}.\over{\unsure{2}}{1}. est donc ramenée à
celle du

\uline{Lemme} \over{\unsure{2}}{3}.\over{\unsure{2}}{1}.1. \inserted{(Zariski)} Soient k un corps, K une extension de type fini de k, \inserted[typed]{m un entier $\geqslant 0$}
$a_i$ $(0\leqslant i\leqslant m)$ des éléments de K \deleted{qui ne sont pas tous algébriques sur} \inserted{tels que le degré de transcendance de $k(a_0,-,a_m)$ sur k soit $\geqslant 2$,}
\deleted{k,} soit $L = K(S_1,..,S_m)$ et k' le sous-corps $k'= k(S_1,..,S_m,a_o +$
$\hand{$\sum_{i}^{m}$}a_iS_i)$ de L (les $S_i$ étant des indéterminées). \deleted{Alors la multiplicité
radicielle de K sur k est égale à celle de L sur k'.}\hand{,} Si K est une
extension primaire de k, alors L est une extension primaire de k'.

\margin{$|$} Ce lemme\hand{,} ou des lemmes qui lui ressemblent comme un frère\hand{,} trai-%
nent un peu partout dans la littérature, c'est pourquoi je te laisse le
choix de l'endroit où tu copieras une démonstration (i.e. je ne me sens
pas inspiré pour en trouver une par mes propres moyens).

\uline{Corollaire} 3.2. Supposons que f soit \deleted{une immersion} \inserted{\unsure{non ramifié}} ou k de car 0,
et que dim f(X)\over{$>$}{$\geqslant$} 2. Alors si X est géom intègre, il en est de même de
Y\hand{$_y$} .

En effet, géom intègre = géom irréd. + séparable.

\uline{Corollaire} 3.3. Supposons \inserted[typed]{k alg. clos, et} que pour toute composante irréductible $X_i$
de X, on ait dim $f(X_i)$\over{$>$}{$\geqslant$} 2, et supposons de plus que X soit \hand{$\unsure{\mathcal{S}}$}-connexe,
où \hand{$\unsure{\mathcal{S}}$} est l'ensemble des parties fermées Z de X telles que dim $\overline{f(\unsure{Z})}=0$\xout{\unsure{x}}
(i.e. pour tout tel Z, X-Z est connexe). Sous ces conditions, Y\hand{$_y$} est
géom connexe sur k(\hand{$y$}).

En effet, par un lemme qui doit figurer au par.6 avec le th. de
Hartshorme, l'hypothèse signifie que l'on peut joindre deux comp. irr.
quelconques X'\xout{\illegible{1}} et X" de X par une chaine de comp irr. $X_o$=X', ...
$X_n$=X", telles que deux consécutives aient une intersection \hand{$\notin\unsure{\mathcal{S}}$}; alors
les images inverses X'\hand{$_y$} et X"\hand{$_y$} sont joints par la chaine des $X_{i}$\hand{$_y$} \xout{\unsure{x}}

\mspage{17}
